\section{Conclusions and Recommendations}

\subsection{Conclusions}
This study successfully developed and evaluated a monocular depth-assisted vision
framework for the non-destructive estimation of the external and internal morphological
properties of Katchamitha (Indian) mangoes in the Philippines. The framework integrated
YOLOv11n instance segmentation with Depth Anything V2 metric depth maps to address the
green-to-green camouflage problem, extracting robust 3D shape and color features from
multi-view images.

Our comparative evaluation of two processing architectures yielded several critical insights:
\begin{enumerate}
    \item \textbf{Instance Segmentation Frontend:} The custom-trained YOLOv11n-seg model
    achieved near-perfect segmentation accuracy (box and mask precision of \textbf{99.9\%},
    recall of \textbf{100.0\%}, and mAP50 score of \textbf{99.5\%}), establishing a robust boundary
    extraction foundation that isolates the fruit from background noise prior to downstream pose
    classification and property estimation.
    \item \textbf{Monocular Scale Ambiguity:} The uncalibrated (No ArUco) pipeline, while
    achieving a robust orientation classification accuracy of \textbf{97.32\%} via the unified
    MobileNetV3-Small CNN classifier and reasonable ripeness classification, failed to
    estimate absolute physical properties (dimensions, volume, and weight). The resulting negative
    or near-zero $R^2$ scores demonstrate that raw monocular depth map features are scale-ambiguous
    and highly sensitive to varying camera distances and lighting, rendering them unsuitable for
    direct physical size estimation without calibration.
    \item \textbf{Fiducial Marker Calibration:} The calibrated (ArUco) pipeline successfully
    resolved the scale ambiguity by using a physical fiducial reference marker to establish an absolute
    scale factor ($S_{\text{aruco}}$). This enabled the conversion of depth coordinates to real-world
    metric values. The calibrated pipeline achieved high morphological estimation accuracy, with
    volumetric $R^2$ scores reaching \textbf{0.83--0.87} and gravimetric $R^2$ scores reaching
    \textbf{0.78--0.86} across different viewpoints, with mass prediction errors kept below
    \textbf{22.8 grams}.
    \item \textbf{Allometric Biological Scaling:} The Layer 2 internal seed regressor achieved
    excellent performance in both configurations ($R^2 = 0.84 - 0.94$). This confirms that the internal
    endocarp's volume and weight scale predictably with the overall size and mass of the external
    fruit, allowing for the non-destructive estimation of the internal stone-to-flesh ratio based on
    external dimensions.
\end{enumerate}

In conclusion, the replacement of the previous feature-based Random Forest orientation
classifiers with a single fine-tuned MobileNetV3-Small CNN improved orientation accuracy
from 75.38\% to \textbf{97.32\%}, demonstrating the value of image-based deep learning
for orientation routing. Furthermore, while monocular depth maps are highly useful for
capturing 3D contours and shape features, absolute metric calibration using a known
physical reference is essential to achieve acceptable accuracy for agricultural sorting and
grading tasks.

\subsection{Recommendations}
Based on the findings of this study, the following recommendations are proposed for future research:
\begin{enumerate}
    \item \textbf{Dataset Expansion and Generalization:} The calibrated pipeline's high
    performance (such as the 100\% ripeness classification accuracy) is constrained by a small
    sample size (10 mangoes with ArUco markers). Future work should expand the calibrated
    dataset to include hundreds of samples. This will provide a more rigorous test support to prevent
    overfitting and confirm the models' generalizability.
    \item \textbf{Reference-Free Metric Scale Estimation:} While the ArUco marker is effective,
    placing a physical marker in the camera's field of view is impractical in automated orchard
    environments. Future research should investigate reference-free depth scaling techniques. This
    includes utilizing smartphone hardware sensors (e.g., dual-camera stereo matching or Time-of-Flight
    LiDAR sensors) or developing deep-learning models trained to infer scale from known agricultural
    features (such as leaf size or trunk thickness).
    \item \textbf{Multi-Variety Validation:} The present study was restricted to Katchamitha
    mangoes. Replicating the dual-pipeline framework on other commercially critical Philippine mango
    varieties, such as the Carabao mango (the country's primary export crop) or the Pico mango,
    will validate the adaptability of the system and support its commercial adoption in precision
    agriculture.
\end{enumerate}
